\documentclass{beamer}
\usepackage{amssymb}
\newcommand{\pair}[2]{(#1,\, #2)}

\begin{document}

  \begin{frame}
    Consideremos la siguiente relación sobre $\mathbb{Z} \times \mathbb{Z}^+$.
  \end{frame}

  \begin{frame}
    Consideremos la siguiente relación sobre $\mathbb{Z} \times \mathbb{Z}^+$.

    \[\pair{x}{y} \equiv \pair{z}{w} \iff x \cdot w = z \cdot y\]
  \end{frame}

  \begin{frame}
    \textit{Teorema}: La relación definida anteriormente es una relación 
    de equivalencia sobre $\mathbb{Z} \times \mathbb{Z}^+$.

  \end{frame}

  \begin{frame}
    \textit{Teorema}: La relación definida anteriormente es una relación 
    de equivalencia sobre $\mathbb{Z} \times \mathbb{Z}^+$.
    \newline\newline

    \textit{Definición}: A la partición asegurada por la relación de equivalencia 
    definida anteriormente, la denotamos $\mathbb{Q} = \{[\pair{x}{y}]: \pair{x}{y} \in \mathbb{Z} \times \mathbb{Z}^+\}$.
  \end{frame}

  \begin{frame}
    Ahora consideramos las siguientes operaciones binarias:
  \end{frame}

  \begin{frame}
    Ahora consideramos las siguientes operaciones binarias:

    \begin{equation*}
      \begin{aligned}
        +: \mathbb{Q} \times \mathbb{Q} &\to \mathbb{Q}: \\
        \pair{\pair{x}{y}}{\pair{w}{z}} &\to \pair{x \cdot z + w \cdot y}{y \cdot z}
      \end{aligned}
    \end{equation*}
  \end{frame}

  \begin{frame}
    Ahora consideramos las siguientes operaciones binarias:

    \begin{equation*}
      \begin{aligned}
        +: \mathbb{Q} \times \mathbb{Q} &\to \mathbb{Q}: \\
        \pair{[\pair{x}{y}]}{[\pair{w}{z}]} &\to [\pair{x \cdot z + w \cdot y}{y \cdot z}]\\
        \cdot : \mathbb{Q} \times \mathbb{Q} &\to \mathbb{Q}: \\
        \pair{[\pair{x}{y}]}{[\pair{w}{z}]} &\to []\pair{x \cdot w}{y \cdot z}]
      \end{aligned}
    \end{equation*}
  \end{frame}

  \begin{frame}
    \textit{Teorema}: Las funciones definidas anteriormente no dependen del elemento 
    escogido.
  \end{frame}

  \begin{frame}
    \textit{Teorema}: Las funciones definidas anteriormente no dependen del elemento 
    escogido.
  \end{frame}

  \begin{frame}
    \textit{Teorema}: Las funciones definidas anteriormente no dependen del elemento 
    escogido.
    \newline\newline

    \textit{Definición}: Denotamos el conjunto de los racionales positivos
    como $\mathbb{Q}^+ = \{\frac{a}{b} \in \mathbb{Q}:  a \in \mathbb{Z}^+\}$.
  \end{frame}

  \begin{frame}
    \textit{Definición}: Definimos la siguiente relación sobre 
    los números racionales:
  \end{frame}

  \begin{frame}
    \textit{Definición}: Definimos la siguiente relación sobre 
    los números racionales:

    \[x \leq y \iff y - x \in \mathbb{Q}^+\]
  \end{frame}

  \begin{frame}
    \textit{Teorema}: Sean $x, y, z, w \in \mathbb{Q}$, entonces se tienen las siguientes propiedades:

    \begin{itemize}
      \item Si $x \leq y$ entonces $x + z \leq x + y$.
    \end{itemize}
  \end{frame}

  \begin{frame}
    \textit{Teorema}: Sean $x, y, z, w \in \mathbb{Q}$, entonces se tienen las siguientes propiedades:

    \begin{itemize}
      \item Si $x \leq y$ entonces $x + z \leq x + y$.
      \item Si $x \leq y$ y $w \leq z$, entonces $x + w \leq y + z$.
    \end{itemize}
  \end{frame}

  \begin{frame}
    \textit{Teorema}: Sean $x, y, z, w \in \mathbb{Q}$, entonces se tienen las siguientes propiedades:

    \begin{itemize}
      \item Si $x \leq y$ entonces $x + z \leq x + y$.
      \item Si $x \leq y$ y $w \leq z$, entonces $x + w \leq y + z$.
      \item Si $x \leq y$ y $z \in \mathbb{Q}^+$, entonces $xz \leq yz$
    \end{itemize}
  \end{frame}

  \begin{frame}
    \textit{Teorema}: Para todo par de números racionales, $x$ e $y$, tal que
    $x < y$, entonces existe un racional $r$ tal que $x < r < y$.
  \end{frame}

  \begin{frame}
    \textit{Teorema}: Para todo par de números racionales, $x$ e $y$, tal que
    $x < y$, entonces existe un racional $r$ tal que $x < r < y$.
    \newline\newline

    \textit{Teorema}: Para todo racional $x \in \mathbb{Q}$ e $y \in \mathbb{Q}^+$, 
    existe un número $n$ entero tal que $x \leq ny$.
  \end{frame}
\end{document}
